
\section{The Free Church’s Mode of Working}

Editorial article in “Fodfehodet” for September 22, 1897.
A rebuttal of J. A. S.’s assertion that the Free Church’s constitution is anarchistic.

The Free Church’s proper fundamental principle is this: that the congregation is the proper divine order within the Christian religion. This principle, which includes within itself that there is no authority which, by ecclesiastical right, may set itself above the congregation or lay claim to its obedience and submission, is, so far as we can see, biblically grounded and justified.

We therefore demand first and foremost of J. A. S. that he must prove from the New Testament that the congregations are obliged to introduce the “representative system,” in order thereby to bring forth a higher unity of the congregations in a church or church body, which is then elevated above them and whose decisions they are obliged to obey, because these church bodies are “authority” over the congregation.

If J. A. S. cannot prove from the New Testament that the congregations are obliged to have synods or associations or annual meetings as authority over themselves and obey their decisions, then all his talk about anarchism and anarchistic principles becomes in reality only a reproach against the Word of God itself, which, so far as we can see, recognizes no authority over the congregations.

It is here that this discussion must seek to find its true ground. It does not help in the least that J. A. S. shows that our constitutional forms of state have introduced a representative system and much compulsion in connection with it.

It does not help a grain. For, in the first place, it is now surely generally acknowledged by all thinking men that there is much imperfection in politics, and that it is not at all a suitable model for the church. Next, there is no state that can drive out the minority or take the right to vote from its representatives, as the United Church, like the Papal Church, is accustomed to do. But finally—and this is absolutely decisive—the Savior himself says definitely and incontrovertibly that “it shall not be so among you.”*

If there is to be any Lutheranism among us, then we must return to the Lutheran principle that Scripture is our only rule and guide. And so long as we, on the ground of Scripture, maintain the congregation’s freedom and authority, then no Lutheran has the right to revile us as anarchists and anarchistic principles.

The example of the state is not binding for the church. The Word of God is our guide. Show us from the Word that the church must work through annual meetings which coerce or expel the congregations, and we will be convinced.

We maintain that according to the Word of God the congregations are fully entitled to decide all their own matters themselves, without an annual meeting or certain association officials deciding them for them.

If now J. A. S. will acknowledge that this our position is the biblical and Lutheran one, then there can further be discussion of how the congregations should cooperate. We must necessarily first get it clarified whether J. A. S. recognizes the decisive authority of Holy Scripture in this question.

Nor is it difficult to show that by the Free Church’s arrangement of purely voluntary meetings, without compulsion upon the congregations, the congregations have a far better opportunity to exercise influence upon the course of affairs than they have in the old associations.

And when J. A. S. means that the same rule that applies to associations must also apply to congregations, then it is precisely this that we deny and reject. The congregation is a body; 3 or 4 congregations, which organize an association, are not a body in the scriptural sense.* They are only bound together by contract, and they are bound only to that which is determined and stipulated in the contract. Why then is not the Free Church’s contract of freedom just as good as the United Church’s contract of compulsion?

Georg Sverdrup